\documentclass[10pt,a4paper]{amsart}
\usepackage{amsmath,amssymb,amsthm,graphicx}
% Retained minimal preamble for the Stacks Project chapter coherent.tex (Stacks chapter, pre-delegation retained fragment for stacks-coherent).
% Verbatim excerpt of the chapter preamble (sources/stacks-spaces-descent/preamble.tex): the reference/slogan/history
% environment block (native lines 7-12), Xy-pic setup (native lines 14-19) and the native theorem declarations with
% the mathematical macros (native lines 154-217). Copied mechanically; nothing added, removed or redefined. The
% wrapper class and the amsmath/amssymb/amsthm/graphicx base packages are supplied by the assembler, so the
% class/xr-hyper/hyperref/externaldocument part of the native preamble is omitted.
% For dealing with references we use the comment environment
\usepackage{verbatim}
\newenvironment{reference}{\comment}{\endcomment}
%\newenvironment{reference}{}{}
\newenvironment{slogan}{\comment}{\endcomment}
\newenvironment{history}{\comment}{\endcomment}
% For commutative diagrams we use Xy-pic
\usepackage[all]{xy}

% We use 2cell for 2-commutative diagrams.
\xyoption{2cell}
\UseAllTwocells
% Theorem environments.
%
\theoremstyle{plain}
\newtheorem{theorem}[subsection]{Theorem}
\newtheorem{proposition}[subsection]{Proposition}
\newtheorem{lemma}[subsection]{Lemma}

\theoremstyle{definition}
\newtheorem{definition}[subsection]{Definition}
\newtheorem{example}[subsection]{Example}
\newtheorem{exercise}[subsection]{Exercise}
\newtheorem{situation}[subsection]{Situation}

\theoremstyle{remark}
\newtheorem{remark}[subsection]{Remark}
\newtheorem{remarks}[subsection]{Remarks}

\numberwithin{equation}{subsection}

% Macros
%
\def\lim{\mathop{\mathrm{lim}}\nolimits}
\def\colim{\mathop{\mathrm{colim}}\nolimits}
\def\Spec{\mathop{\mathrm{Spec}}}
\def\Hom{\mathop{\mathrm{Hom}}\nolimits}
\def\Ext{\mathop{\mathrm{Ext}}\nolimits}
\def\SheafHom{\mathop{\mathcal{H}\!\mathit{om}}\nolimits}
\def\SheafExt{\mathop{\mathcal{E}\!\mathit{xt}}\nolimits}
\def\Sch{\mathit{Sch}}
\def\Mor{\mathop{\mathrm{Mor}}\nolimits}
\def\Ob{\mathop{\mathrm{Ob}}\nolimits}
\def\Sh{\mathop{\mathit{Sh}}\nolimits}
\def\NL{\mathop{N\!L}\nolimits}
\def\CH{\mathop{\mathrm{CH}}\nolimits}
\def\proetale{{pro\text{-}\acute{e}tale}}
\def\etale{{\acute{e}tale}}
\def\QCoh{\mathit{QCoh}}
\def\Ker{\mathop{\mathrm{Ker}}}
\def\Im{\mathop{\mathrm{Im}}}
\def\Coker{\mathop{\mathrm{Coker}}}
\def\Coim{\mathop{\mathrm{Coim}}}

% Boxtimes
%
\DeclareMathSymbol{\boxtimes}{\mathbin}{AMSa}{"02}

%
% Macros for moduli stacks/spaces
%
\def\QCohstack{\mathcal{QC}\!\mathit{oh}}
\def\Cohstack{\mathcal{C}\!\mathit{oh}}
\def\Spacesstack{\mathcal{S}\!\mathit{paces}}
\def\Quotfunctor{\mathrm{Quot}}
\def\Hilbfunctor{\mathrm{Hilb}}
\def\Curvesstack{\mathcal{C}\!\mathit{urves}}
\def\Polarizedstack{\mathcal{P}\!\mathit{olarized}}
\def\Complexesstack{\mathcal{C}\!\mathit{omplexes}}
% \Pic is the operator that assigns to X its picard group, usage \Pic(X)
% \Picardstack_{X/B} denotes the Picard stack of X over B
% \Picardfunctor_{X/B} denotes the Picard functor of X over B
\def\Pic{\mathop{\mathrm{Pic}}\nolimits}
\def\Picardstack{\mathcal{P}\!\mathit{ic}}
\def\Picardfunctor{\mathrm{Pic}}
\def\Deformationcategory{\mathcal{D}\!\mathit{ef}}
\usepackage{cleveref}
\crefname{originallemma}{original lemma}{original lemmas}
\Crefname{originallemma}{Original lemma}{Original lemmas}
\makeatletter
\newlabel{cohomology-lemma-Leray}{{01F2}{source}{}{sourcepointer-1}{}}
\newlabel{cohomology-lemma-Leray@cref}{{[originallemma][0][]01F2}{[1][0][]source}{}{}{}}
\makeatother
\crefname{originallemma}{original lemma}{original lemmas}
\Crefname{originallemma}{Original lemma}{Original lemmas}
\makeatletter
\newlabel{cohomology-section-locality}{{01E0}{source}{}{sourcepointer-2}{}}
\newlabel{cohomology-section-locality@cref}{{[originallemma][0][]01E0}{[1][0][]source}{}{}{}}
\makeatother
\crefname{originallemma}{original lemma}{original lemmas}
\Crefname{originallemma}{Original lemma}{Original lemmas}
\makeatletter
\newlabel{lemma-quasi-coherent-affine-cohomology-zero}{{01XB}{source}{}{sourcepointer-3}{}}
\newlabel{lemma-quasi-coherent-affine-cohomology-zero@cref}{{[originallemma][0][]01XB}{[1][0][]source}{}{}{}}
\makeatother
\crefname{originallemma}{original lemma}{original lemmas}
\Crefname{originallemma}{Original lemma}{Original lemmas}
\makeatletter
\newlabel{lemma-quasi-compact-h1-zero-covering}{{01XF}{source}{}{sourcepointer-4}{}}
\newlabel{lemma-quasi-compact-h1-zero-covering@cref}{{[originallemma][0][]01XF}{[1][0][]source}{}{}{}}
\makeatother
\crefname{originallemma}{original lemma}{original lemmas}
\Crefname{originallemma}{Original lemma}{Original lemmas}
\makeatletter
\newlabel{properties-lemma-extend-trivial}{{01PE}{source}{}{sourcepointer-5}{}}
\newlabel{properties-lemma-extend-trivial@cref}{{[originallemma][0][]01PE}{[1][0][]source}{}{}{}}
\makeatother
\crefname{originallemma}{original lemma}{original lemmas}
\Crefname{originallemma}{Original lemma}{Original lemmas}
\makeatletter
\newlabel{schemes-lemma-push-forward-quasi-coherent}{{01LC}{source}{}{sourcepointer-6}{}}
\newlabel{schemes-lemma-push-forward-quasi-coherent@cref}{{[originallemma][0][]01LC}{[1][0][]source}{}{}{}}
\makeatother
\crefname{originallemma}{original lemma}{original lemmas}
\Crefname{originallemma}{Original lemma}{Original lemmas}
\makeatletter
\newlabel{schemes-lemma-quasi-compact-permanence}{{03GI}{source}{}{sourcepointer-7}{}}
\newlabel{schemes-lemma-quasi-compact-permanence@cref}{{[originallemma][0][]03GI}{[1][0][]source}{}{}{}}
\makeatother
\crefname{originallemma}{original lemma}{original lemmas}
\Crefname{originallemma}{Original lemma}{Original lemmas}
\makeatletter
\newlabel{schemes-lemma-quasi-compact-preserved-base-change}{{01K5}{source}{}{sourcepointer-8}{}}
\newlabel{schemes-lemma-quasi-compact-preserved-base-change@cref}{{[originallemma][0][]01K5}{[1][0][]source}{}{}{}}
\makeatother
\title{Relative affineness criterion via first cohomology of ideal sheaves: a quasi-compact morphism f : X -> Y with X, Y quasi-separated is affine as soon as R\textasciicircum{}1f\_*I vanishes for every quasi-coherent sheaf of ideals on X.}
\author{The Stacks Project Authors}
\date{stacks-project:coherent.tex:local-snapshot-20261012; GNU FDL 1.2 or later}
\begin{document}
\maketitle
\noindent\textit{Editorial note.} Statement span is lines 718-723 of sources/stacks-coherent/pinned-original-source.tex; proof span is lines 725-746, the complete author proof (20 body lines).\par\medskip
\noindent\textit{Editorial note.} Navigation stubs: 8 cited-but-not-included labels are supplied as non-load-bearing pointers with their pinned tags where the snapshot tag index has them.\par\medskip
\noindent\textit{Editorial note.} extra\_packages is empty on purpose: the retained preamble already loads xy with options and a second declaration clashes in the assembler wrapper.\par\medskip
\noindent\textit{Editorial note.} Text is retained verbatim, including the author's own typos and placeholders.\par\medskip
\noindent\textit{Editorial note.} Pool stacks-coherent backfill wave 33e.\par\medskip
\noindent\textit{Editorial note.} Feeds lemma-characterize-finite (card L16381) which uses the same criterion family; shelf chain consistent.\par\medskip
\noindent\textit{Allowed standard prerequisites (editorial).}\begin{itemize}
\item Explicit prerequisite boundary. See why-field; every pointer named here occurs verbatim in the retained statement/proof text.
\end{itemize}\medskip
\ifdefined\hypertarget\hypertarget{sourcepointer-1}{}\fi
\noindent\textit{Source reference.} Original lemma 01F2: sources/stacks-coherent/pinned-original-source.tex, source lines 746--746. Cited by the retained statement or proof but not part of the collected spans; supplied only as a navigation pointer. The argument does not depend on its content.\par\medskip
\ifdefined\hypertarget\hypertarget{sourcepointer-2}{}\fi
\noindent\textit{Source reference.} Original lemma 01E0: sources/stacks-coherent/pinned-original-source.tex, source lines 746--746. Cited by the retained statement or proof but not part of the collected spans; supplied only as a navigation pointer. The argument does not depend on its content.\par\medskip
\ifdefined\hypertarget\hypertarget{sourcepointer-3}{}\fi
\noindent\textit{Source reference.} Original lemma 01XB: sources/stacks-coherent/pinned-original-source.tex, source lines 363--367. Cited by the retained statement or proof but not part of the collected spans; supplied only as a navigation pointer. The argument does not depend on its content.\par\medskip
\ifdefined\hypertarget\hypertarget{sourcepointer-4}{}\fi
\noindent\textit{Source reference.} Original lemma 01XF: sources/stacks-coherent/pinned-original-source.tex, source lines 500--503. Cited by the retained statement or proof but not part of the collected spans; supplied only as a navigation pointer. The argument does not depend on its content.\par\medskip
\ifdefined\hypertarget\hypertarget{sourcepointer-5}{}\fi
\noindent\textit{Source reference.} Original lemma 01PE: sources/stacks-coherent/pinned-original-source.tex, source lines 746--746. Cited by the retained statement or proof but not part of the collected spans; supplied only as a navigation pointer. The argument does not depend on its content.\par\medskip
\ifdefined\hypertarget\hypertarget{sourcepointer-6}{}\fi
\noindent\textit{Source reference.} Original lemma 01LC: sources/stacks-coherent/pinned-original-source.tex, source lines 746--746. Cited by the retained statement or proof but not part of the collected spans; supplied only as a navigation pointer. The argument does not depend on its content.\par\medskip
\ifdefined\hypertarget\hypertarget{sourcepointer-7}{}\fi
\noindent\textit{Source reference.} Original lemma 03GI: sources/stacks-coherent/pinned-original-source.tex, source lines 746--746. Cited by the retained statement or proof but not part of the collected spans; supplied only as a navigation pointer. The argument does not depend on its content.\par\medskip
\ifdefined\hypertarget\hypertarget{sourcepointer-8}{}\fi
\noindent\textit{Source reference.} Original lemma 01K5: sources/stacks-coherent/pinned-original-source.tex, source lines 746--746. Cited by the retained statement or proof but not part of the collected spans; supplied only as a navigation pointer. The argument does not depend on its content.\par\medskip
\begin{lemma}
\label{lemma-criterion-affine-morphism}
Let $f : X \to Y$ be a quasi-compact morphism with $X$ and $Y$ quasi-separated.
If $R^1f_*\mathcal{I} = 0$ for every quasi-coherent sheaf of ideals
$\mathcal{I}$ on $X$, then $f$ is affine.
\end{lemma}

\begin{proof}
Let $V \subset Y$ be an affine open subscheme. We have to show that
$U = f^{-1}(V)$ is affine. The inclusion morphism $V \to Y$ is quasi-compact
by Schemes, Lemma \ref{schemes-lemma-quasi-compact-permanence}.
Hence the base change $U \to X$ is quasi-compact, see
Schemes, Lemma \ref{schemes-lemma-quasi-compact-preserved-base-change}.
Thus any quasi-coherent sheaf of ideals $\mathcal{I}$ on $U$
extends to a quasi-coherent sheaf of ideals on $X$, see
Properties, Lemma \ref{properties-lemma-extend-trivial}.
Since the formation of $R^1f_*$ is local on $Y$
(Cohomology, Section \ref{cohomology-section-locality})
we conclude that $R^1(U \to V)_*\mathcal{I} = 0$ by the assumption
in the lemma. Hence by the Leray Spectral sequence
(Cohomology, Lemma \ref{cohomology-lemma-Leray})
we conclude that $H^1(U, \mathcal{I}) = H^1(V, (U \to V)_*\mathcal{I})$.
Since $(U \to V)_*\mathcal{I}$ is quasi-coherent by
Schemes, Lemma \ref{schemes-lemma-push-forward-quasi-coherent}, we have
$H^1(V, (U \to V)_*\mathcal{I}) = 0$ by
Lemma \ref{lemma-quasi-coherent-affine-cohomology-zero}.
Thus we find that $U$ is affine by
Lemma \ref{lemma-quasi-compact-h1-zero-covering}.
\end{proof}

\end{document}
